% Main matter control is handled by main.tex; remove pandoc-specific includes.

\hypertarget{introduction}{%
\section{Introduction}\label{introduction}}

\hypertarget{statement-of-purpose-objectives}{%
\subsection{Statement of purpose:
objectives}\label{statement-of-purpose-objectives}}

This thesis is one of two coordinated master's theses on the same
project. The companion thesis {[}Motje, 2026{]} covers the
\textbf{software} side (gateway firmware, services, authentication,
monitoring, and the application stack). The present document covers the
\textbf{hardware} side (the network fabric and the endpoint fleet) and
the \textbf{knowledge artefact} that ties both halves together.

The five objectives below were defined jointly with the companion
thesis. They are common to both documents; what differs is the dimension
of each objective that each author addresses.

\begin{enumerate}
\def\labelenumi{\arabic{enumi}.}
\tightlist
\item
  \textbf{Create a living document} that any contributor with basic
  technical literacy can update, so that the knowledge it contains grows
  with the projects that produce it.
\item
  \textbf{Develop a high-level guide} for readers with basic technical
  knowledge --- not a reference manual aimed at networking specialists.
\item
  \textbf{Make the guide useful} for organisations such as AUCOOP, whose
  volunteers rotate frequently and whose deployments are geographically
  distributed.
\item
  \textbf{Validate the guide} by applying it in a real deployment:
  specifically, the Namaqua Kalahari Children Hope project at N
  Mutschuana Primary School, Gochas, Namibia.
\item
  \textbf{Centralise}, in a single openly maintained artefact, the
  knowledge accumulated from previous AUCOOP projects, so that future
  cohorts inherit it rather than reconstruct it.
\end{enumerate}

The present document covers the \textbf{hardware dimension} of these
objectives: the chapters of the handbook that deal with the wireless
network fabric (planning, IP addressing, mesh backhaul, antennas, power,
and enclosure) and with the endpoint pipeline (refurbished laptop
intake, AUCOOP golden-master image, and PXE/Clonezilla mass deployment).
It also reports the validation of those chapters through the network and
endpoint deployment in Namibia (March 2026).

\hypertarget{motivation}{%
\subsection{1.2 Motivation}\label{motivation}}

The motivation for this thesis comes from two distinct but reinforcing
sources. The first is \emph{external}: the persistent gap in basic
connectivity and usable ICT equipment in the communities where AUCOOP
(\emph{Associació d'Universitaris per la Cooperació}) can make a
difference. The second is \emph{internal}: a recurring pain point in the
association, namely the loss of knowledge between volunteer cohorts.
This loss forces members to re-solve the same problems with each new
generation of volunteers.

Either of the two would justify a thesis on its own. Together they
explain why the deliverable of this work is not a single deployment but
a \emph{reusable instrument} --- the handbook described in
\autoref{sec:handbook} --- built and validated through a real field
deployment.

\hypertarget{external-motivation-connectivity-and-usable-endpoints-as-preconditions-for-everything-else}{%
\subsubsection{External motivation --- connectivity and usable endpoints
as preconditions for everything
else}\label{external-motivation-connectivity-and-usable-endpoints-as-preconditions-for-everything-else}}

Reliable connectivity is no longer a luxury for the communities AUCOOP
partners with: it is a transformative tool that creates new
opportunities for individuals. Schools, health posts, community centres,
and NGOs increasingly depend on online services for curricula,
record-keeping, and communication with regional authorities. In the
contexts AUCOOP has worked in over the last decade --- rural Senegal,
the Namibian Hardap region, and several pilots inside Catalonia --- the
obstacle is rarely a single missing piece. It is the combination of
three deficits that reinforce each other:

\begin{enumerate}
\def\labelenumi{\arabic{enumi}.}
\tightlist
\item
  \textbf{No usable network.} Either there is no infrastructure at all,
  or there is a single ADSL/4G uplink at one location, with no capacity
  to reach the classrooms that need it.
\item
  \textbf{No usable endpoints.} Even when connectivity is solved, the
  community often has only a handful of working computers, all aged,
  often with broken storage or pirated operating systems that cannot be
  patched. New hardware is unaffordable. Donations of refurbished
  hardware exist, but they require expertise to receive, image, and
  deploy at scale.
\item
  \textbf{No usable documentation.} The reference material that does
  exist assumes a level of bandwidth, budget, and expertise that the
  field site does not have. Vendor manuals presume current firmware and
  online activation. Academic papers describe ideal architectures but
  skip the thirty practical decisions a deployment team faces in an
  afternoon.
\end{enumerate}

This thesis attacks the first two deficits directly --- through the
network work in \autoref{sec:methodology-network} and the endpoint
reconditioning pipeline in \autoref{sec:methodology-endpoint}. It
attacks the third deficit by producing the handbook described in
\autoref{sec:handbook}: a genuine, opinionated, field-tested alternative
to the documentation that is currently missing.

\hypertarget{internal-motivation-knowledge-continuity-at-aucoop}{%
\subsubsection{Internal motivation --- knowledge continuity at
AUCOOP}\label{internal-motivation-knowledge-continuity-at-aucoop}}

\href{https://aucoop.upc.edu}{AUCOOP} is a student volunteer association
at UPC. Its strength is also its structural weakness: it is staffed by
bachelor and master students, who join enthusiastic, contribute for one
to three academic years, graduate, and leave. Over more than three years
inside the association, the author has observed the same anti-pattern
repeatedly:

\begin{itemize}
\tightlist
\item
  A project ends. Its deliverable --- a working network, a configured
  server, a refurbished classroom --- exists physically, but its
  \emph{knowledge} exists only in the heads of the students who built it
  and in scattered private notes (chat messages, personal Markdown
  files, hand-drawn diagrams in field notebooks).
\item
  Those students graduate within twelve months. The institutional memory
  of the project leaves with them.
\item
  The next cohort, picking up either a follow-up project at the same
  site or a similar project elsewhere, finds something in the shared
  drive --- but no clear way to apply it. They start from a blank page,
  rediscover the same OpenWrt mesh quirks, run into the same DHCP
  option, hit the same Clonezilla
  \texttt{partclone\ target\ seek\ ERROR}, and pay the same debugging
  cost the previous cohort already paid.
\item
  Three years later, the cycle repeats.
\end{itemize}

The cause is not laziness; it is the absence of a \emph{living artefact}
in which contribution is the natural side-effect of doing the work. PDFs
in a shared drive are not such an artefact: they are written
\emph{after} the project, by people already mentally checked out, and
read by nobody. Internal wikis on self-hosted servers are not such an
artefact either: they go unmaintained the moment the volunteer who set
them up leaves, and they are invisible to external partners.

The thesis takes the position that this internal problem is solvable,
and that solving it is itself a worthwhile engineering contribution. The
chosen solution --- an open Git repository of Markdown sources that
builds to a public website \emph{and} a downloadable PDF, governed by
lightweight contribution rules and assisted by AI-driven authoring
tooling --- is described in \autoref{sec:handbook}. It is already in
use: this thesis was written against it, the Namibia 2026 deployment was
documented in it, and the next AUCOOP student to pick up a project will
find a non-empty starting point.

\hypertarget{personal-motivation}{%
\subsubsection{Personal motivation}\label{personal-motivation}}

A short note on personal motivation. The author has been involved with
AUCOOP since the bachelor years and has been on the receiving end of
both deficits described above: arriving at a project with no
documentation, and leaving a project knowing that the documentation she
wrote would probably not be read. This thesis is the deliberate attempt
to break that pattern on the way out --- to leave behind something that
the next person can stand on, rather than start beside.

\hypertarget{requirements-and-specifications}{%
\subsection{Requirements and
specifications}\label{requirements-and-specifications}}

The handbook --- the deliverable that ties together the network and
endpoint contributions described in this thesis --- must satisfy five
requirements. These follow from the problem statement of
\autoref{sec:motivation} and the objectives of \autoref{sec:objectives}.
They are common to both this thesis and the companion software thesis
{[}Motje, 2026{]}, and are reproduced here in the formulation jointly
agreed with that work, adapted where the hardware-side framing differs.

\textbf{R1 --- Living and contributor-friendly.} The artefact must be
\emph{easier to update than to discard}. This rules out the conventional
NGO-report format --- a monolithic PDF authored once by a departing
volunteer and then frozen --- as the primary form of the handbook. It
points instead to a source-text format under version control (Markdown
in a public Git repository), with a low barrier to contribution: short
pull requests, lightweight review, and contribution rules small enough
to read in one sitting. This requirement is the direct technical answer
to the volunteer-rotation problem of \autoref{sec:motivation-internal}:
if updating the artefact is more expensive than reinventing the
knowledge it contains, no rational volunteer will update it.

\textbf{R2 --- Openly accessible.} The handbook must be reachable from
the field with intermittent connectivity, and consultable when there is
no connectivity at all. Public web hosting alone is therefore
insufficient: an exportable form (a downloadable full PDF, mirrored on
the same site) is required, so that a volunteer arriving in Gochas or in
rural Senegal can consult the same material on a laptop that last saw a
network three weeks ago. The dual-output build pipeline
(\autoref{sec:handbook} and \autoref{appendix:online-handbook})
implements this requirement.

\textbf{R3 --- Pedagogically progressive.} The handbook must serve
readers ranging from those without prior exposure to community networks
through to volunteers needing a recipe for a specific task on a specific
morning. A single linear document cannot serve both audiences without
alienating one of them. The handbook addresses this through two
complementary tracks: a narrative \emph{Imaginary Use Case} (Chapter 2
of the handbook) that introduces the domain through a story, and a
topical \emph{Guide} (Chapter 3) organised as recipes. The two tracks
are maintained in strict 1:1 correspondence, so that the same
operational content is reachable both narratively and topically. The
structural choice is described in \autoref{sec:handbook-structure}.

\textbf{R4 --- Free and open-source software (FOSS) where the recipe
controls the stack.} The recommendations the handbook makes must be
reproducible without commercial licences and maintainable by future
cohorts using only freely available tooling. For the software the recipe
configures --- operating systems, network firmware, services,
documentation toolchain --- this requirement is strict: OpenWrt rather
than vendor firmware, Linux Mint rather than Windows, ISC \texttt{dhcpd}
rather than a commercial appliance, OnlyOffice Community rather than a
licensed suite. For the hardware the recipe \emph{deploys onto}, the
requirement is necessarily looser: vendor BIOS firmware on Lenovo
ThinkPads and the binary radio firmware blobs that OpenWrt loads onto
Cudy access points are accepted as out-of-scope dependencies, because no
FOSS alternative exists for either at this date, and refusing the
donation pipeline on those grounds would defeat the larger refurbishment
goal. The requirement is therefore \emph{FOSS for everything we author
or configure; vendor firmware tolerated where unavoidable, and recorded
explicitly in the relevant recipe.}

\textbf{R5 --- Field-validated.} The recommendations the handbook
contains must be traceable to a real deployment, not invented from
documentation alone. Every recipe in Chapter 3 of the handbook that this
thesis covers --- wireless mesh, IP addressing, network planning, laptop
deployment, AUCOOP image --- is paired with an entry in Chapter 4 of the
handbook that records where it was first exercised in the field. The
validation strategy is summarised in \autoref{sec:results-strategy}, and
the eleven operational lessons that emerged from the Namibia deployment
are consolidated in \autoref{sec:lessons-consolidated}.

A sixth, more conventional set of hardware-recipe requirements applies
to the deployments themselves rather than to the handbook as artefact:

\begin{longtable}[]{@{}
  >{\raggedright\arraybackslash}p{(\columnwidth - 4\tabcolsep) * \real{0.3333}}
  >{\raggedright\arraybackslash}p{(\columnwidth - 4\tabcolsep) * \real{0.3333}}
  >{\raggedright\arraybackslash}p{(\columnwidth - 4\tabcolsep) * \real{0.3333}}@{}}
\toprule\noalign{}
\begin{minipage}[b]{\linewidth}\raggedright
\#
\end{minipage} & \begin{minipage}[b]{\linewidth}\raggedright
Hardware-recipe requirement
\end{minipage} & \begin{minipage}[b]{\linewidth}\raggedright
Why
\end{minipage} \\
\midrule\noalign{}
\endhead
\bottomrule\noalign{}
\endlastfoot
H1 & Network gear must be supported by OpenWrt or an equivalent FOSS
firmware & Auditability, longevity beyond vendor support window \\
H2 & Endpoints must be sourced from refurbishment, not new manufacture &
\autoref{sec:embedded-carbon-avoided} (carbon) and
\autoref{sec:budget-comparison} (cost) \\
H3 & Mass-deployment time per endpoint must be field-viable & The
deployment window is one-two work days \\
H4 & The total per-site cash BOM must be raisable from a single small
NGO grant & Replicability inside AUCOOP and similar associations
(cf.~\autoref{sec:budget-funding}) \\
\end{longtable}

H1--H4 are operationalised in the recipes; their validation against the
Namibia deployment is reported in \autoref{sec:results-network} and
\autoref{sec:results-endpoint}.

\hypertarget{methods-and-procedures-relation-to-prior-and-concurrent-work}{%
\subsection{Methods and procedures --- relation to prior and concurrent
work}\label{methods-and-procedures-relation-to-prior-and-concurrent-work}}

This thesis is \textbf{methodologically integrated} with two adjacent
bodies of work, and explicitly delimited from a third. Naming each is
part of the contribution, because the value of the thesis lies as much
in \emph{what it does not re-do} as in what it adds.

\textbf{Prior work --- the Hahatay deployment in Senegal.} The work
presented here is, in its origin, a continuation of the deployment that
an AUCOOP team --- including the author --- carried out in Hahatay
(Senegal) over more than three years of sustained engagement. Hahatay is
one of the flagship projects of the association: a multi-site community
network built incrementally, without a structured guide, in a rural
Senegalese context against constraints recognisably similar to those of
the Namibia case (intermittent connectivity, refurbished endpoints,
volunteer rotation between cohorts). The Hahatay deployment is the
principal source of the operational knowledge captured in the Community
Network Handbook. The Namibia deployment described in
\autoref{ch:results} is therefore not the first exercise of the recipes
in the handbook, but the first \emph{field validation of the recipes
after they had been written down}: the Hahatay engagement produced the
knowledge, the handbook codified it, and Namibia tested whether the
codification travelled.

\textbf{Prior work --- the AUCOOP Community Network Handbook itself.}
The handbook exists as a public Git repository {[}AUCOOP, 2026{]} and
was seeded with the Hahatay knowledge described above before this thesis
began. The thesis treats the handbook as both input and output: input,
because much of the network-planning, IP-addressing, and OpenWrt
material was authored before the Namibia deployment and was used to plan
it; output, because every operational lesson from Namibia was committed
back into the handbook in real time.

\textbf{Concurrent work --- the companion software thesis.} The
companion thesis {[}Motje, 2026{]} covers the services layer (Proxmox
virtualisation, Nextcloud collaborative platform, RADIUS authentication,
captive portal, Zabbix monitoring, VPN backhaul, and backup procedures).
The two theses share a common Chapter 1 (introduction), Chapter 2 (state
of the art), and Chapter 3 (methodology) in their respective
formulations; they share the Namibia deployment as a common validation
site; and they share a bibliography backbone (entries marked
\texttt{{[}shared{]}} in \autoref{ch:bibliography}). They diverge in the
technical contribution: the present document does not describe the
services running on top of the network, and the companion thesis does
not describe the network or the endpoints they run on. Where the
boundary is operationally fuzzy --- DHCP, DNS, basic monitoring --- this
thesis treats those services as boundary infrastructure required for the
hardware to function, and describes only what is necessary for that
purpose; the deeper services treatment is in {[}Motje, 2026{]}.

\textbf{External tooling reused, not implemented.} This work uses
extensively the open-source ecosystem developed by the wider
community-network movement. None of the components below are
contributions of this thesis; the contribution lies in the way they are
\textbf{integrated, configured, and explained} in the context of small
community deployments executed by volunteer teams.

\begin{itemize}
\tightlist
\item
  \textbf{Network firmware and stack.} OpenWrt 24.10.x as the
  access-point and gateway operating system; the \texttt{mac80211}
  802.11s mesh stack; \texttt{hostapd}, \texttt{wpa\_supplicant},
  \texttt{dnsmasq}, ISC \texttt{dhcpd}, and \texttt{tftpd-hpa} for
  boundary services.
\item
  \textbf{Endpoint provisioning.} Clonezilla and \texttt{partclone} for
  bit-level imaging and restore; GRUB 2 and \texttt{pxelinux} for the
  PXE boot chain; the NFS kernel server for image distribution.
\item
  \textbf{Endpoint operating system.} Linux Mint Cinnamon 22.3 as the
  base of the AUCOOP golden-master image; OnlyOffice Community as the
  bundled productivity suite.
\item
  \textbf{Refurbishment provenance.} The eReuse / DeviceHub toolchain
  for tracking donated hardware from intake to deployed.
\item
  \textbf{Documentation toolchain.} MkDocs with the Material/Zensical
  themes for building the handbook to both web and PDF; GitHub Actions
  for continuous integration; a public GitHub repository for hosting and
  contribution governance.
\end{itemize}

The author's contribution lies in the \textbf{integration},
\textbf{field validation}, \textbf{operational distillation}, and
\textbf{knowledge-artefact} layers; co-authorship of individual handbook
chapters is recorded per-commit in the handbook git history.

\textbf{Out of scope.} Custom firmware development, novel routing
protocols, formal performance modelling of mesh throughput, and
academic-grade life-cycle assessment of refurbishment are all out of
scope. The thesis is an integration and validation effort, not a
research contribution to any of those sub-fields.
\autoref{sec:env-limits} makes the LCA limitation explicit;
\autoref{sec:future-work} lists the others as future-work directions.

\hypertarget{work-plan-and-gantt}{%
\subsection{Work plan and Gantt}\label{work-plan-and-gantt}}

See \protect\hyperlink{16-gantt-diagram}{\texttt{gantt}} below.

\hypertarget{phases}{%
\subsubsection{Phases}\label{phases}}

\begin{longtable}[]{@{}
  >{\raggedright\arraybackslash}p{(\columnwidth - 4\tabcolsep) * \real{0.3333}}
  >{\raggedright\arraybackslash}p{(\columnwidth - 4\tabcolsep) * \real{0.3333}}
  >{\raggedright\arraybackslash}p{(\columnwidth - 4\tabcolsep) * \real{0.3333}}@{}}
\toprule\noalign{}
\begin{minipage}[b]{\linewidth}\raggedright
Phase
\end{minipage} & \begin{minipage}[b]{\linewidth}\raggedright
Description
\end{minipage} & \begin{minipage}[b]{\linewidth}\raggedright
Period
\end{minipage} \\
\midrule\noalign{}
\endhead
\bottomrule\noalign{}
\endlastfoot
P1 & Bootstrap handbook repository, conventions, CI & Sep 2025 -- Oct
2025 \\
P2 & Network chapters (planning, IP, mesh, antennas, power) & Oct 2025
-- Jan 2026 \\
P3 & Laptop chapters (refurbishment, AUCOOP image, PXE/Clonezilla) & Jan
2026 -- Mar 2026 \\
P4 & Field deployment --- Namibia (Gochas) & Mar 2026 \\
P5 & Real-use case write-up + handbook consolidation & Mar 2026 -- Apr
2026 \\
P6 & Thesis write-up & Apr 2026 -- May 2026 \\
P7 & Defence preparation & May 2026 \\
\end{longtable}

\hypertarget{gantt-diagram}{%
\subsection{Gantt diagram}\label{gantt-diagram}}

\begin{quote}
Source: \texttt{plan/gantt.md}. Render via
\texttt{latex/gantt\_diagrama.tex} (uses \texttt{pgfgantt}).
\end{quote}

\begin{figure}[ht]
  \centering
  % Gantt chart — Project phases P1–P7 (Sep 2025 – May 2026)
%
% Grid: 18 half-month slots, 2 per calendar month.
%   Slot  1– 2  → Sep 2025   (1=early Sep, 2=late Sep)
%   Slot  3– 4  → Oct 2025
%   Slot  5– 6  → Nov 2025
%   Slot  7– 8  → Dec 2025
%   Slot  9–10  → Jan 2026
%   Slot 11–12  → Feb 2026
%   Slot 13–14  → Mar 2026   (13=early Mar ≈ 1–15, 14=late Mar ≈ 16–31)
%   Slot 15–16  → Apr 2026   (15=early Apr ≈ 1–15, 16=late Apr ≈ 16–30)
%   Slot 17–18  → May 2026
%
% P4 field trip  20 Mar – 14 Apr  →  slots 14–15  (late-Mar … early-Apr)
%
% Included via % Gantt chart — Project phases P1–P7 (Sep 2025 – May 2026)
%
% Grid: 18 half-month slots, 2 per calendar month.
%   Slot  1– 2  → Sep 2025   (1=early Sep, 2=late Sep)
%   Slot  3– 4  → Oct 2025
%   Slot  5– 6  → Nov 2025
%   Slot  7– 8  → Dec 2025
%   Slot  9–10  → Jan 2026
%   Slot 11–12  → Feb 2026
%   Slot 13–14  → Mar 2026   (13=early Mar ≈ 1–15, 14=late Mar ≈ 16–31)
%   Slot 15–16  → Apr 2026   (15=early Apr ≈ 1–15, 16=late Apr ≈ 16–30)
%   Slot 17–18  → May 2026
%
% P4 field trip  20 Mar – 14 Apr  →  slots 14–15  (late-Mar … early-Apr)
%
% Included via % Gantt chart — Project phases P1–P7 (Sep 2025 – May 2026)
%
% Grid: 18 half-month slots, 2 per calendar month.
%   Slot  1– 2  → Sep 2025   (1=early Sep, 2=late Sep)
%   Slot  3– 4  → Oct 2025
%   Slot  5– 6  → Nov 2025
%   Slot  7– 8  → Dec 2025
%   Slot  9–10  → Jan 2026
%   Slot 11–12  → Feb 2026
%   Slot 13–14  → Mar 2026   (13=early Mar ≈ 1–15, 14=late Mar ≈ 16–31)
%   Slot 15–16  → Apr 2026   (15=early Apr ≈ 1–15, 16=late Apr ≈ 16–30)
%   Slot 17–18  → May 2026
%
% P4 field trip  20 Mar – 14 Apr  →  slots 14–15  (late-Mar … early-Apr)
%
% Included via \input{latex/gantt_diagrama.tex} inside a figure environment.
\begin{ganttchart}[
  hgrid,
  vgrid={
    *{1}{draw=gray!25,thin},          % mid-month divider (odd slots)
    *{1}{draw=gray!60,semithick},     % month boundary  (even slots)
    *{1}{draw=gray!25,thin},
    *{1}{draw=gray!60,semithick},
    *{1}{draw=gray!25,thin},
    *{1}{draw=gray!60,semithick},
    *{1}{draw=gray!25,thin},
    *{1}{draw=gray!60,semithick},
    *{1}{draw=gray!25,thin},
    *{1}{draw=gray!60,semithick},
    *{1}{draw=gray!25,thin},
    *{1}{draw=gray!60,semithick},
    *{1}{draw=gray!25,thin},
    *{1}{draw=gray!60,semithick},
    *{1}{draw=gray!25,thin},
    *{1}{draw=gray!60,semithick},
    *{1}{draw=gray!25,thin},
    *{1}{draw=gray!60,semithick}
  },
  x unit=0.68cm,
  y unit title=0.55cm,
  y unit chart=0.65cm,
  title height=1,
  title label font=\bfseries\footnotesize,
  bar label font=\small,
  bar label node/.append style={align=right, text width=5.2cm},
  bar/.style={draw=black!70, fill=cyan!55},
  bar height=0.55,
  milestone/.append style={fill=red!70, draw=red!80, shape=diamond},
  milestone height=0.45,
]{1}{18}

  %% ---- Row 1: month labels (span 2 slots each) ---------------------------
  \gantttitle{Sep}{2}
  \gantttitle{Oct}{2}
  \gantttitle{Nov}{2}
  \gantttitle{Dec}{2}
  \gantttitle{Jan}{2}
  \gantttitle{Feb}{2}
  \gantttitle{Mar}{2}
  \gantttitle{Apr}{2}
  \gantttitle{May}{2} \\

  %% ---- Row 2: year markers ------------------------------------------------
  \gantttitle{2025}{8}
  \gantttitle{2026}{10} \\

  %% ---- Phases -------------------------------------------------------------
  % P1  Sep 2025 – Oct 2025   → slots 1–4
  \ganttbar[name=P1]{P1 — Handbook bootstrap \& CI}{1}{4} \\

  % P2  Oct 2025 – Jan 2026   → slots 3–10
  \ganttbar[name=P2]{P2 — Network chapters}{3}{10} \\

  % P3  Jan 2026 – Mar 2026   → slots 9–14
  \ganttbar[name=P3]{P3 — Laptop chapters}{9}{14} \\

  % P4  20 Mar – 14 Apr 2026  → slots 14–15  (late-Mar … early-Apr)
  \ganttbar[name=P4,
    bar/.style={draw=black!70, fill=orange!70}
  ]{P4 — Field deployment (Namibia)}{14}{15} \\

  % P5  Mar 2026 – Apr 2026   → slots 13–16
  \ganttbar[name=P5]{P5 — Write-up \& consolidation}{13}{16} \\

  % P6  Apr 2026 – May 2026   → slots 15–18
  \ganttbar[name=P6]{P6 — Thesis write-up}{15}{18} \\

  % P7  May 2026              → slots 17–18
  \ganttbar[name=P7,
    bar/.style={draw=black!70, fill=green!50}
  ]{P7 — Defence preparation}{17}{18} \\

  %% ---- Submission milestone (end of May) ----------------------------------
  \ganttmilestone[name=M4]{Submission}{18}

\end{ganttchart}
 inside a figure environment.
\begin{ganttchart}[
  hgrid,
  vgrid={
    *{1}{draw=gray!25,thin},          % mid-month divider (odd slots)
    *{1}{draw=gray!60,semithick},     % month boundary  (even slots)
    *{1}{draw=gray!25,thin},
    *{1}{draw=gray!60,semithick},
    *{1}{draw=gray!25,thin},
    *{1}{draw=gray!60,semithick},
    *{1}{draw=gray!25,thin},
    *{1}{draw=gray!60,semithick},
    *{1}{draw=gray!25,thin},
    *{1}{draw=gray!60,semithick},
    *{1}{draw=gray!25,thin},
    *{1}{draw=gray!60,semithick},
    *{1}{draw=gray!25,thin},
    *{1}{draw=gray!60,semithick},
    *{1}{draw=gray!25,thin},
    *{1}{draw=gray!60,semithick},
    *{1}{draw=gray!25,thin},
    *{1}{draw=gray!60,semithick}
  },
  x unit=0.68cm,
  y unit title=0.55cm,
  y unit chart=0.65cm,
  title height=1,
  title label font=\bfseries\footnotesize,
  bar label font=\small,
  bar label node/.append style={align=right, text width=5.2cm},
  bar/.style={draw=black!70, fill=cyan!55},
  bar height=0.55,
  milestone/.append style={fill=red!70, draw=red!80, shape=diamond},
  milestone height=0.45,
]{1}{18}

  %% ---- Row 1: month labels (span 2 slots each) ---------------------------
  \gantttitle{Sep}{2}
  \gantttitle{Oct}{2}
  \gantttitle{Nov}{2}
  \gantttitle{Dec}{2}
  \gantttitle{Jan}{2}
  \gantttitle{Feb}{2}
  \gantttitle{Mar}{2}
  \gantttitle{Apr}{2}
  \gantttitle{May}{2} \\

  %% ---- Row 2: year markers ------------------------------------------------
  \gantttitle{2025}{8}
  \gantttitle{2026}{10} \\

  %% ---- Phases -------------------------------------------------------------
  % P1  Sep 2025 – Oct 2025   → slots 1–4
  \ganttbar[name=P1]{P1 — Handbook bootstrap \& CI}{1}{4} \\

  % P2  Oct 2025 – Jan 2026   → slots 3–10
  \ganttbar[name=P2]{P2 — Network chapters}{3}{10} \\

  % P3  Jan 2026 – Mar 2026   → slots 9–14
  \ganttbar[name=P3]{P3 — Laptop chapters}{9}{14} \\

  % P4  20 Mar – 14 Apr 2026  → slots 14–15  (late-Mar … early-Apr)
  \ganttbar[name=P4,
    bar/.style={draw=black!70, fill=orange!70}
  ]{P4 — Field deployment (Namibia)}{14}{15} \\

  % P5  Mar 2026 – Apr 2026   → slots 13–16
  \ganttbar[name=P5]{P5 — Write-up \& consolidation}{13}{16} \\

  % P6  Apr 2026 – May 2026   → slots 15–18
  \ganttbar[name=P6]{P6 — Thesis write-up}{15}{18} \\

  % P7  May 2026              → slots 17–18
  \ganttbar[name=P7,
    bar/.style={draw=black!70, fill=green!50}
  ]{P7 — Defence preparation}{17}{18} \\

  %% ---- Submission milestone (end of May) ----------------------------------
  \ganttmilestone[name=M4]{Submission}{18}

\end{ganttchart}
 inside a figure environment.
\begin{ganttchart}[
  hgrid,
  vgrid={
    *{1}{draw=gray!25,thin},          % mid-month divider (odd slots)
    *{1}{draw=gray!60,semithick},     % month boundary  (even slots)
    *{1}{draw=gray!25,thin},
    *{1}{draw=gray!60,semithick},
    *{1}{draw=gray!25,thin},
    *{1}{draw=gray!60,semithick},
    *{1}{draw=gray!25,thin},
    *{1}{draw=gray!60,semithick},
    *{1}{draw=gray!25,thin},
    *{1}{draw=gray!60,semithick},
    *{1}{draw=gray!25,thin},
    *{1}{draw=gray!60,semithick},
    *{1}{draw=gray!25,thin},
    *{1}{draw=gray!60,semithick},
    *{1}{draw=gray!25,thin},
    *{1}{draw=gray!60,semithick},
    *{1}{draw=gray!25,thin},
    *{1}{draw=gray!60,semithick}
  },
  x unit=0.68cm,
  y unit title=0.55cm,
  y unit chart=0.65cm,
  title height=1,
  title label font=\bfseries\footnotesize,
  bar label font=\small,
  bar label node/.append style={align=right, text width=5.2cm},
  bar/.style={draw=black!70, fill=cyan!55},
  bar height=0.55,
  milestone/.append style={fill=red!70, draw=red!80, shape=diamond},
  milestone height=0.45,
]{1}{18}

  %% ---- Row 1: month labels (span 2 slots each) ---------------------------
  \gantttitle{Sep}{2}
  \gantttitle{Oct}{2}
  \gantttitle{Nov}{2}
  \gantttitle{Dec}{2}
  \gantttitle{Jan}{2}
  \gantttitle{Feb}{2}
  \gantttitle{Mar}{2}
  \gantttitle{Apr}{2}
  \gantttitle{May}{2} \\

  %% ---- Row 2: year markers ------------------------------------------------
  \gantttitle{2025}{8}
  \gantttitle{2026}{10} \\

  %% ---- Phases -------------------------------------------------------------
  % P1  Sep 2025 – Oct 2025   → slots 1–4
  \ganttbar[name=P1]{P1 — Handbook bootstrap \& CI}{1}{4} \\

  % P2  Oct 2025 – Jan 2026   → slots 3–10
  \ganttbar[name=P2]{P2 — Network chapters}{3}{10} \\

  % P3  Jan 2026 – Mar 2026   → slots 9–14
  \ganttbar[name=P3]{P3 — Laptop chapters}{9}{14} \\

  % P4  20 Mar – 14 Apr 2026  → slots 14–15  (late-Mar … early-Apr)
  \ganttbar[name=P4,
    bar/.style={draw=black!70, fill=orange!70}
  ]{P4 — Field deployment (Namibia)}{14}{15} \\

  % P5  Mar 2026 – Apr 2026   → slots 13–16
  \ganttbar[name=P5]{P5 — Write-up \& consolidation}{13}{16} \\

  % P6  Apr 2026 – May 2026   → slots 15–18
  \ganttbar[name=P6]{P6 — Thesis write-up}{15}{18} \\

  % P7  May 2026              → slots 17–18
  \ganttbar[name=P7,
    bar/.style={draw=black!70, fill=green!50}
  ]{P7 — Defence preparation}{17}{18} \\

  %% ---- Submission milestone (end of May) ----------------------------------
  \ganttmilestone[name=M4]{Submission}{18}

\end{ganttchart}

  \caption{Project phases and schedule (Sep 2025 — May 2026)}
\end{figure}

\hypertarget{document-structure}{%
\subsection{Document structure}\label{document-structure}}

The remainder of this document is organised as follows.

\autoref{ch:state-of-the-art} reviews the state of the art across the
areas relevant to the work: existing community networks (Guifi.net, NYC
Mesh, Freifunk, AlterMundi, Rhizomatica), the ICT-for-development and
refurbishment ecosystem (Labdoo, eReuse, the WEEE regulatory frame), the
OpenWrt firmware ecosystem and the open-source mass-deployment toolchain
(Clonezilla, PXE), and the literature on knowledge management in
volunteer organisations (tacit and explicit knowledge, working
knowledge, communities of practice, the volunteer-rotation problem).

\autoref{ch:methodology} describes the methodology in three parts that
correspond to the three contributions of the thesis:
\autoref{sec:methodology-network} the network hardware deployment,
organised through a four-functional-layer model (field, site, endpoint
touchpoint, power and enclosure); \autoref{sec:methodology-endpoint} the
endpoint reconditioning and mass-deployment pipeline (intake, AUCOOP
image, PXE/Clonezilla); and \autoref{sec:handbook} the AUCOOP Community
Network Handbook itself as a knowledge artefact, including the choice of
documentation tooling, the Diátaxis-inspired Ch2/Ch3 split, the
contribution governance, and the AI-assisted authoring workflow.

\autoref{ch:results} reports the results of applying the methodology at
N Mutschuana Primary School, Gochas (Namibia), in March - April 2026.
The validation is organised through three qualitative analytical lenses
--- \emph{coverage}, \emph{sufficiency}, \emph{adaptation} --- and the
operational lessons it produced are consolidated in
\autoref{sec:lessons-consolidated}, together with the handbook artefact
each one fed back into.

\autoref{ch:budget} presents the budget, distinguishing the cash outlay
an organisation actually has to raise from the full project cost at
indicative engineering rates, and comparing both with the
commercial-equivalent alternative.

\autoref{ch:environmental-impact} addresses environmental impact, with
particular attention to the manufacturing-phase carbon dominance that
justifies the refurbishment requirement, and to the limits of the
analysis (this thesis does not constitute a formal life-cycle
assessment).

\autoref{ch:conclusions} draws conclusions against each of the five
objectives stated in \autoref{sec:objectives}, records the limitations
honestly alongside the conclusions they qualify, lists future-work
directions, and documents the hand-over protocol for the next AUCOOP
cohort.

The bibliography (\autoref{ch:bibliography}) marks shared entries with
the companion thesis explicitly. \autoref{appendix:online-handbook}
points to the publicly hosted handbook (URL, build instructions,
contribution conventions); \autoref{appendix:glossary} is the glossary
of acronyms; and the remaining appendices collect the budget detail (BOM
and itemised costs), the deployment field log, representative
configuration excerpts, and the data sheet for the Namibia case.

\hypertarget{deviations-from-the-initial-plan}{%
\subsection{Deviations from the initial
plan}\label{deviations-from-the-initial-plan}}

The plan submitted at the start of the project (September 2025) held in
its broad outline but required four significant adjustments during
execution. Recording them here explains schedule shifts and the
practical choices the team made.

\begin{enumerate}
\def\labelenumi{\arabic{enumi}.}
\tightlist
\item
  Routers acquired at the beginning changed hardware and were no longer
  OpenWrt-compatible. We replaced them with Cudy devices that had
  reliable OpenWrt support.
\item
  The field work was planned for January 2026, but the Children's House
  inauguration that month and availability constraints from the local
  counterpart forced the team to move the deployment to late March /
  early April 2026.
\item
  The donated computers shipped with a preinstalled operating system
  that users found unfriendly. Based on their feedback, the project
  scoped and produced an AUCOOP golden-master image (Linux Mint based)
  and added a PXE-based deployment workflow to provision the +10
  computers with the open-source, user-friendly image.
\item
  Deployment at the Children's House was not possible because the
  ministry had not provided the ADSL router. We prepared the mesh there
  in advance, but the live network could not be completed until the
  uplink was available.
\end{enumerate}

Each of these deviations is discussed where relevant in
\autoref{sec:methodology-network} and
\autoref{sec:methodology-endpoint}.
